\providecommand{\simroot}{.}
\section{Supplementary Numerical Results}
\label{app:simulation_results}

\subsection{Additional Channel-Bank Comparisons}
The reduced-data MNIST comparisons use 10,000 training examples, a hidden
width of 128, and smaller analog-channel banks. They provide a sensitivity
check rather than an independent test dataset: the official test set is
the same as in the full-data comparison. Table~\ref{tab:app_reduced_mnist}
reports all seeds of the corresponding ten-draw cohorts. A separate
five-draw screening cohort is retained in
Table~\ref{tab:app_screen_mnist}; its numbers are not pooled with the
ten-draw evaluations.
\input{\simroot/tables/appendix_reduced_mnist}
\input{\simroot/tables/appendix_screen_mnist}

The width-64 feature-erasure refinements in
Table~\ref{tab:app_toy_erasure} use the same deployment probability $0.5$
but different bank sizes. They show the performance attainable with a
larger transmitted representation. The proposed variants do not improve
on ERM at every bank size. This complements the main finite-channel
mechanism, rather than establishing uniform superiority across channels.
\input{\simroot/tables/appendix_toy_erasure}

\subsection{Data Size and Optimization Budget}
Table~\ref{tab:app_data_size} varies the number of training examples at
$m=100$. Accuracy increases with data size in these runs, but the protocol
also increases the number of minibatch updates because it fixes epochs,
not updates. The training/validation realizations change with $n$, and the
listed repetitions vary initialization and channel banks rather than
resampling independent datasets. Moreover, the observed signed data gap
is negative in these realizations. This experiment therefore describes
the data-size performance trend; it does not independently verify the
theoretical $n$-rate.
\input{\simroot/tables/appendix_data_size}

\subsection{Internal Objective Ablations and Selected Cases}
The proposed objectives themselves belong to the PAC-Bayesian regularizer
family. The data-only penalty is consequently an internal ablation. Both
T1 terms depend on the same posterior KL; when that KL dominates the
confidence logarithms, the separate penalties are approximately a scaled
square-root KL penalty. T2 has a single square-root KL form with adjusted
scaling and confidence. An unrestricted retuning of a data-only
coefficient can therefore yield near-equivalent optimization. The
distinct theoretical contribution is the treatment of finite channel
observations and its resulting dependence on $n$ and $m$.

In the full-data comparison, the validation-selected T1 and data-only
configurations attain $76.6307\%$ and $76.6370\%$, respectively, a
practical tie. Allowing test-based hyperparameter selection over the
entire trained grid instead gives the data-only control $76.7513\%$.
Completing the missing Fashion-MNIST learning-rate/noise controls reduces
the corresponding 100-draw T1 advantage to only $0.0208$ percentage point.
These findings limit claims of a fundamentally different regularizer.

Selected favorable checkpoints remain useful as illustrative cases. For
full-data MNIST at training seed 9001, a new 100-draw stream gives T1
$75.9712\%$, exceeding the strongest of the three reevaluated data-only
coefficients at $75.7801\%$. For the reduced-data $m=50$ case, seed 1003
gives T2 $67.7729\%$ versus $67.4309\%$. These are explicitly
test/seed-selected examples. Their Monte Carlo intervals condition on
the fixed checkpoints and test data, and do not establish a population
advantage across training seeds or undo selection effects. They are not
used in the main mean-performance table.

\subsection{Selected Optimization Settings}
Learning rates in the grayscale-image searches are $0.001$ and $0.003$, with
posterior standard-deviation initializations $0.003$ and $0.01$. The
initial L2 search uses coefficients $10^{-5}$ and $10^{-4}$ coupled to
these two learning rates. The grid is therefore not a full factorial or
exhaustive L2 optimization. The initial T1 coefficient pairs are
$(\eta_1,\eta_2)=(a,10a)$ with
$a\in\{0.0003,0.001,0.003,0.01\}$, and T2 uses joint coefficients $2a$.
Learning rate and initial posterior noise are conditionally assigned,
rather than searched as a full Cartesian product. Perfect-channel
posterior training considers data coefficients $0.006$ and $0.06$.
Full-data refinement adds T1 pairs $(0.001,0.01)$ and $(0.003,0.03)$,
and T2 coefficients $0.0015$ and $0.004$, at learning rate $0.003$ and
initial posterior standard deviation $0.01$. The exact selected values
are listed in Table~\ref{tab:app_hyperparameters}. All weights and biases are
included in the explicit squared-weight L2 penalty. The complete
conditional search configurations are retained in the accompanying
reproducibility manifests.
\input{\simroot/tables/selected_hyperparameters}

\subsection{CIFAR-10 Screening and Optimization Settings}
The preliminary CIFAR-10 screen uses 10,000 training examples, 5,000
validation examples, 30 epochs, and two initialization/channel-bank seeds.
Table~\ref{tab:app_cifar_screen} retains both seeds of every selected
configuration and uses ten fresh deployment draws. T1 and T2 exceed ERM
and L2 in both seeds, but ERM has lower bounded log loss. The accuracy and
bounded-risk orderings therefore differ in this screen.
\input{\simroot/tables/appendix_cifar_screen}

The 34 training trials cover learning rates $0.001$ and $0.003$ for ERM;
both learning rates crossed with L2 coefficients $10^{-6}$ and $10^{-5}$;
and both learning rates crossed with the two proposed settings
$(s_0,a)=(0.001,0.0003)$ and $(0.003,0.001)$, where $s_0$ is the initial
posterior standard deviation. T1 uses $(\eta_1,\eta_2)=(a,10a)$ and T2
uses joint coefficient $2a$. Each perfect-training reference has one
screened configuration: learning rate $0.001$, L2 coefficient $10^{-6}$
where applicable, and initial posterior standard deviation $0.001$ with
data coefficient $0.01$ for Proposed Perfect. These references are not
exhaustively optimized over perfect-training hyperparameters.

The full-data extension freezes the seven validation-selected settings
and uses seeds 5101--5103 with 45,000 training examples and 100 epochs.
ERM retains learning rate $0.001$; matched L2, T1 and T2 retain $0.003$.
Matched L2 selects $10^{-6}$; T1 selects $(0.001,0.01)$ and T2 selects
joint coefficient $0.002$, both with $s_0=0.003$. The two stages share
the validation prefix and a seeded data permutation, with the smaller
training subset nested inside the larger one. They are separate
data-size/optimization-budget stages on the same official test set.
